\section{Contraction over coordinate sweeps}\label{sweeps:section}
The block moments and merge identity of
Section~\ref{cycles:block-core} apply to every deterministic signed
input. Here we sharpen their two-sweep bounds and then compute the
entire one-sweep quadratic form by a collision coupling. The latter
retains the densities in the sister blocks met by a pair of walkers.
For a random initial list those densities are balanced at each fixed
time, yielding a different averaged contraction. We also record the
layer-by-layer alternative supplied by refreshed coordinate halves.

\subsection{Energy and sequential revelation}\label{sweeps:mass-section}
We use the signed-weight process of Section~\ref{cycles:weights}, with
the coordinate order $1,\ldots,d$. To distinguish a realized energy
from its expectation, write $H_t$ for the available set, $x_t$ for the
signed vector, and $E_t=\|x_t\|_2^2$. For a tagged card after preceding
paths are exposed,
\[
 x_t=p_t-|H_t|^{-1}\one_{H_t}.
\]
This is the vector $f_t$ of Section~\ref{sec:first-route}. Its conditional
law $p_t$ is defined given the exposed paths, but its value at time $t$
is computed from their prefixes. All expectations below average those
paths under the original shuffle law. We also use deterministic
zero-sum signed inputs, for which $E_t\le E_0$ pathwise.

\begin{lemma}[Sequential comparison]\label{sweeps:sequential}
Suppose that for every preceding set in an ordered list of $k$ cards, the
corresponding tagged process satisfies $\E E_T\le a$. Then the joint
endpoint law of the list has total-variation distance at most
$k\sqrt{na}/2$ from the uniform injection. The same conclusion holds after
averaging over a random ordered list if the energy bounds hold in that average.
\end{lemma}
\begin{proof}
The conditional one-card error is
$\E\|x_T\|_1/2\le\sqrt{|H_T|\E E_T}/2\le\sqrt{na}/2$.
Lemma~\ref{lem:sequential-tv} sums these errors, first coarsening the
preceding paths to their endpoints. Averaging the resulting inequality
over the initial list gives the last assertion.
\end{proof}

\subsection{Completed-sweep symmetry and explicit bounds}\label{sweeps:blocks}
Retain the level-$j$ blocks $\mathcal B_j$ of
Section~\ref{cycles:block-core}. For the output of a first sweep, put
$M_B=\sum_{u\in B}x(u)$ and $h_B=|H\cap B|$. Thus $M_B,h_B$ are
the quantities denoted $S_B,m_B$ in
Lemma~\ref{sweeps:block-statistics}. Before applying their moment bounds
again, we give a second proof that identifies a symmetry of the whole
output law. This symmetry will be strengthened in
Section~\ref{prefix:section}.

\begin{proof}[A symmetry proof of Lemma~\ref{sweeps:block-statistics}]
Fix $B\in\mathcal B_j$, put $s=2^j$, and let $m$ be the conserved
number of available positions. Immediately after updating coordinate $t$, the law is invariant under swapping
any chosen set of its matching edges: edges with one available position have independent symmetric
outputs and outputs on edges with two available positions are equal.  A symmetry that preserves the later
matchings remains a symmetry through the later updates.  In particular, after
a sweep its law is invariant under independently translating the suffix
$(j+1,\ldots,d)$ at each fixed length-$j$ prefix.  To verify this assertion,
flip one coordinate $t>j$ on all sites with a given prefix.  This is a union of
coordinate-$t$ edges, and it preserves every later matching.  Such flips
generate all the asserted suffix translations.

Apply independent uniform suffix translations to an already completed
configuration.  Conditional on that configuration, the entries selected for
$B$ sample one site uniformly and independently from each prefix fiber.  Their
signed sum has mean zero and variance at most $sE_0/n$; their available count is a
sum of independent indicators with mean $ms/n$.  The translated configuration
has the original law, so the same variance and exponential-moment estimates
prove the lemma.  This argument establishes the needed transitivity without
asserting independence among the entries of the original completed sweep.
\end{proof}

\begin{proposition}[Uniform two-sweep bound]\label{sweeps:two-sweep}
Fix $0<\varepsilon\le1$ and put $h=\lfloor d/2\rfloor$,
$\sigma=1-\varepsilon/4$.  For deterministic zero-sum signed mass supported
on at least $\varepsilon n$ available positions, two sweeps satisfy
\begin{equation}\label{sweeps:two-sweep-explicit}
 \mathbb E E_{2d}\le E_0\left[
 \sigma^{d-h}+2^{-h}+(d+1)n e^{-\varepsilon^2 2^h/2}\right].
\end{equation}
In particular, for every
$0<b<\min\{-\tfrac12\log_2(1-\varepsilon/4),\tfrac12\}$,
$\mathbb E E_{2d}\le n^{-b}E_0$ for all sufficiently large $d$.
For $\varepsilon=1/16$, one may take $b=1/200$.
\end{proposition}
\begin{proof}
Call the first-sweep output good if every block of level at least $h$ has
available-position density at least $\varepsilon/2$.  Lemma~\ref{sweeps:block-statistics}
and a union bound give the displayed exceptional probability.  Conditional
on a good output, let $F_j$ be the expected energy after $j$ layers of the
second sweep.  The merging identity and $2M_1M_2\le M_1^2+M_2^2$ give
\[
 F_{j+1}\le\sigma F_j+2^{-j-1}\sum_{B\in\mathcal B_j}M_B^2,
 \qquad h\le j<d.
\]
We have $F_h\le E_0$ and
$\mathbb E\sum_{B\in\mathcal B_j}M_B^2\le E_0$ by
\eqref{sweeps:block-variance}.  Iterate the inequality, bound its geometric
weights by one, and sum $\sum_{j=h}^{d-1}2^{-j-1}\le2^{-h}$.
On the exceptional event use $E_{2d}\le E_0$.  This proves
\eqref{sweeps:two-sweep-explicit}.  Its first term has exponent
$-\tfrac12\log_2\sigma$, its second has exponent $1/2$, and the final term
decays faster than every power of $n$.  For $\varepsilon=1/16$,
$-\tfrac12\log_2(63/64)>1/200$, proving the last assertion.
\end{proof}

Two variants of the displayed bound retain useful quantitative information.
The symmetry proof and the count of fewer than $2n$ blocks replace the last
term by $2n\exp[-c\varepsilon2^h]$, with
$c=1/2-e^{-1}$.  For $\varepsilon=1/8$, starting the second-sweep recursion
at $\lceil d/2\rceil$ gives the more explicit estimate
\begin{equation}\label{sweeps:eighth-bound}
 \frac{\mathbb E E_{2d}}{E_0}
 \le(31/32)^{\lfloor d/2\rfloor}
      +16n^{-1/2}+32n e^{-\sqrt n/64}
 \le n^{-1/100}
\end{equation}
for sufficiently large $d$ (the ratio is interpreted only when $E_0>0$;
when $E_0=0$ all energies vanish).  Indeed a fixed block of size $s$ has
failure probability at most $e^{-s/64}$ by the Bernoulli lower-tail bound,
and the per-layer errors are $E_0/(2s)+nE_0e^{-s/64}$.
Their geometric sum is bounded using $(1-31/32)^{-1}=32$.

Because the two-sweep estimate is uniform in the deterministic input, it can
be iterated conditionally.  With $\varepsilon=1/16$ and $r=2048$, after
$2r$ sweeps the tag energy is at most $n^{-r/200}$ in expectation.  Hence
Lemma~\ref{sweeps:sequential} gives uniformly, for every list of at most
$15n/16$ cards, endpoint distance at most
\begin{equation}\label{sweeps:colored-marginal}
 \tfrac12 n^{3/2-r/400}=o(1).
\end{equation}
Similarly, \eqref{sweeps:eighth-bound} gives the uniform $7n/8$-card marginal
estimate after $2r$ sweeps for any fixed integer $r>300$.

\subsection{An exact collision formula}\label{sweeps:collision-section}
The preceding proof bounds a merge directly.  The next calculation instead
computes the complete one-sweep quadratic form by coupling two conditional
walkers.  It retains more information about the initial available-position configuration.
For $u\in V$, let $C_j(u)$ be its level-$j$ block.  For an available set $H$ and
$B\in\mathcal B_i$, define
\[
 s_H(B)=\prod_{j=i+1}^d(1-\eta_j(B)/2),\qquad
 \eta_j(B)=\frac{|H\cap(C_j(u)\setminus C_{j-1}(u))|}{2^{j-1}}
 \quad(u\in B).
\]
The blocks in this formula are independent of the choice of $u\in B$.
For a singleton $\{u\}$ use $i=0$, and an empty product is one.

\begin{lemma}[One-sweep collision identity]\label{sweeps:collision}
For deterministic $(H,x)$ at the start of a sweep,
\begin{align}\label{sweeps:collision-eq}
 \mathbb E E_d={}&\sum_{u\in H}s_H(\{u\})x(u)^2\\
 &+\sum_{i=1}^d2^{1-i}\sum_{B\in\mathcal B_i}
          s_H(B)x(B_0)x(B_1),\nonumber
\end{align}
where $x(A)=\sum_{u\in A}x(u)$ and $B_0,B_1$ are the children of $B$.
\end{lemma}
\begin{proof}
Given the complete exposed paths, let $M(u,v)$ be the conditional
transition probability for a hidden tag from $u$ to $v$. Linearity of
the signed flow gives
\[
 x_d(v)=\sum_{u\in H}x(u)M(u,v).
\]
Consequently the coefficient
$\E\sum_vM(u,v)M(z,v)$ in its expected squared norm is a collision
probability: sample the exposed paths with their actual law, and then
sample two walkers independently conditional on those paths, starting
at $u,z\in H$.

We analyze this joint law chronologically. Generate the physical
switches and the resulting available sets in layer order. On a pair
with one available input, both walkers use its forced available
output if they are there. On a pair with two available inputs, each
walker uses its own extra fair coin, independent of every physical
switch and of the other walker's coins. Conditional on the complete
exposed paths, this construction gives exactly two independent copies
of $M$, by Lemma~\ref{lem:marginal-averaging}. In the following
calculation we average over the paths; we do not retain full-path
conditioning when invoking independence of physical switch arrays.

For $u\ne z$, let $i$ be their largest differing coordinate, and for
$u=z$ put $i=0$. Before coordinate $i$, they use distinct matching
pairs. Each updated bit is fair, and the two bits are independent
given the preceding chronological history. If their first $i-1$
outputs agree, their coordinate-$i$ inputs form a pair of available
sites. Their independent auxiliary choices then agree with probability
$1/2$. Their probability of agreement through coordinate $i$ is
therefore $2^{-i}$.

Now take $j>i$. Reveal the physical switches in the block
$C_{j-1}(u)$ during its first $j-1$ layers and the auxiliary choices
of both walkers through that time. Their histories are confined to
this block, so this information determines whether they agree and,
on agreement, their common address. No switches in its sister block
$C_j(u)\setminus C_{j-1}(u)$ have been revealed. Those switches are
independent of the revealed information. Each site of the sister
block has expected availability equal to its initial density
$\eta_j(C_i(u))$, since its first $j-1$ coordinate averages have
all been applied. This holds in particular at the now determined
partner address.

If that partner is unavailable, the common input has a forced output
and the walkers stay together. If it is available, their independent
auxiliary choices agree with probability $1/2$. Hence each level
$j>i$ contributes the factor $1-\eta_j(C_i(u))/2$. Any disagreement
in a processed coordinate is permanent within the sweep. Multiplying
the successive conditional probabilities gives
\[
 \E\sum_vM(u,v)M(z,v)=2^{-i}s_H(C_i(u)),
\]
including $i=0$. Group the ordered off-diagonal pairs by their largest
differing coordinate. The two orders of each pair give the coefficient
$2^{1-i}$ in \eqref{sweeps:collision-eq}.
\end{proof}

\begin{proposition}[Averaged two-sweep contraction]\label{sweeps:averaged}
Choose the initial ordered card list uniformly and independently of all
switches.  Fix an exposure index with $m\ge\varepsilon n$ available positions, where
$\varepsilon=1/q$ and $q$ is a fixed power of two.  Put
$b=\lfloor d/2\rfloor$ and
\[
 \Delta_n=2n(n+1)\exp[-\varepsilon(1/2-e^{-1})2^b].
\]
At sweep boundaries $t$, the tag process obeys
\begin{equation}\label{sweeps:averaged-recursion}
 \mathbb E E_{t+2d}\le\Delta_n+
 \left[(1+b/2)(1-\varepsilon/4)^{d-b}+2^{-b}\right]\mathbb E E_t.
\end{equation}
With $q=128$, $R=128q$, and $T=2Rd$,
$\mathbb E E_T\le n^{-8}+R\Delta_n$.  Consequently the average endpoint
TV distance of an ordered list of $n-n/q$ cards tends to zero.
\end{proposition}
\begin{proof}
At every fixed time the available set is a uniform $m$-subset: after fixing all
switches the shuffle is a permutation independent of the random initial list.
This is an unconditional statement.  Call an available set good if every block
of level at least $b$ has density at least $\varepsilon/2$.  It suffices by
coupling to consider $m=\varepsilon n$.  A Bernoulli-$\varepsilon$ subset
conditioned on size $\varepsilon n$ is uniform, and this conditioning event
has probability at least $1/(n+1)$ because $\varepsilon n$ is a mode of the
binomial distribution.  The exponential-moment estimate in
Lemma~\ref{sweeps:block-statistics}, followed by a union bound over fewer
than $2n$ blocks, gives $\Pr(H_t\text{ not good})\le\Delta_n$.

Put $w=x_{t+d}$.  On a good second-sweep input,
$s_{H_{t+d}}(B)\le(1-\varepsilon/4)^{d-b}$ for every level at most $b$.
In \eqref{sweeps:collision-eq} use $2w(B_0)w(B_1)\le w(B_0)^2+w(B_1)^2$.
At levels $1\le i\le b$,
$\sum_{C\in\mathcal B_{i-1}}w(C)^2\le2^{i-1}\|w\|_2^2$;
each such level contributes at most half the common factor times energy.
The remaining levels contribute at most
$\sum_{i=b+1}^d2^{-i}\sum_{C\in\mathcal B_{i-1}}w(C)^2$.
On failure of goodness, use $E_{t+2d}\le1$.  Apply the block-variance bound
conditionally over the first sweep to each last sum; after summing over a
partition its expectation is at most $\mathbb E E_t$.  This proves
\eqref{sweeps:averaged-recursion} without assuming that energy and goodness
are independent.

For $\gamma=\varepsilon/16$, the bracket is at most $n^{-\gamma}$ for all
sufficiently large $d$: its first term is at most
$(1+b/2)e^{-\varepsilon d/8}$ and its second is at most $2n^{-1/2}$.
Iteration through $R$ sweep pairs gives
$\mathbb E E_T\le n^{-R\gamma}+R\Delta_n=n^{-8}+R\Delta_n$.
The last term decreases faster than every power of $n$.
Lemma~\ref{sweeps:sequential}, averaged over lists, now proves the assertion.
\end{proof}

\subsection{Refreshing coordinate halves}\label{sweeps:refresh-section}
A separate argument contracts energy after each layer following an initial
$2d$ delay.  It is useful when one wants estimates uniform over the initial
list without first grouping layers into two-sweep blocks.
Write $i_t=1+(t\bmod d)$ and use the coordinate averages $P_i$
of Section~\ref{sec:first-route}.
Let $\mathcal F_t$ contain all switch outcomes strictly before layer $t$.
The path-defined signed-mass process is adapted to this larger filtration.
Its defining conditional card law still conditions only on exposed paths,
not on all variables in $\mathcal F_t$.

\begin{lemma}[Noise and refreshed balance]\label{sweeps:refresh}
For the tagged-card difference $x_t=p_t-m^{-1}{\bf1}_{H_t}$ with a fixed
initial hidden set of size $m\ge\varepsilon n$, so that $E_0\le1$, put
$a_t=\mathbb E E_t$ and
$I_t=\sum_u x_t(u)^2{\bf1}_{\{u\oplus e_{i_t}\notin H_t\}}$.
Then
\begin{equation}\label{sweeps:noise-identity}
 a_t=\sum_{j=1}^d2^{-j}\mathbb E I_{t-j}\quad(t\ge d).
\end{equation}
For $t\ge d$, each of the following choices gives
$\mathbb E I_t\le\gamma a_t+B_n$:
\[
 \begin{array}{c|c|c}
 &\gamma&B_n\\ \hline
 \text{available count}&1-\varepsilon/2&2e^{-\varepsilon n/8}\\
 \text{blocker count}&1-\varepsilon/2&2e^{-\varepsilon^2n/4}.
 \end{array}
\]
For either choice, every $1/2<\beta<1$ with $2\beta-1\ge\gamma$ satisfies
\begin{equation}\label{sweeps:refresh-bound}
 a_t\le\beta^{t-2d}+B_n/(1-\gamma)\qquad(t\ge0).
\end{equation}
\end{lemma}
\begin{proof}
Apply Lemma~\ref{noise:memory} with $f_t=x_t$, $V_t=H_t$, and
$w_t=\E I_t$ to obtain \eqref{sweeps:noise-identity}. The opposite-half
factorization of Lemma~\ref{mart:halves} holds also with the larger
past field $\mathcal F_{t-d+1}$: after that time the two halves use
independent fresh switch arrays, and the path-defined vector in either
half is a function of its own array. It remains to supply a lower
bound on both available densities at that conditioning time.

The available-count argument in the proof of Lemma~\ref{mart:balance}
bounds the chance that either half has fewer than $m/4$ available
positions by $2e^{-m/8}\le2e^{-\varepsilon n/8}$. Off that event,
both densities are at least $\varepsilon/2$. The factorization and
$E_t\le1$ on the exception give the first row. For the second row,
Lemma~\ref{noise:halves} with $\alpha=\varepsilon$ gives
$\E I_t\le(1-\varepsilon/2)a_t+2e^{-\varepsilon^2n/4}E_0$;
use $E_0\le1$.

Finally use \eqref{noise:comparison-solution} with
$\rho=\gamma$, $\eta=B_n$, $t_0=2d$, and decay parameter $\beta$.
Its condition $\rho/(2\beta-1)\le1$ is exactly the stated hypothesis.
\end{proof}

For the blocker-count form, retain the explicit choices
$\varepsilon=1/32$, $\beta=1-\varepsilon/8$, and $T=2048d$.
Then $\beta^{T-2d}\le n^{-5}$, and
Lemma~\ref{sweeps:sequential} gives $o(1)$ uniformly for all lists of at most
$31n/32$ cards.  For the available-count form, take $\varepsilon=1/128$,
$\beta=1-\varepsilon/4$, and any integer $b_0>2$ with
$\beta^{b_0-2}\le2^{-6}$.  At $T=b_0d$ the per-tag expected TV error is
$o(n^{-1})$, hence every list of $127n/128$ cards has $o(1)$ endpoint error.
The tail estimates above give both refreshed-balance bounds.


\subsection{Two ways to dominate complementary cosets}\label{sweeps:truncation}
Let $G=S_n$ and partition the input positions into $q$ equal sets
$A_1,\ldots,A_q$.  Let $K_i=\operatorname{Sym}(A_i)$ fix the complement
pointwise.  The coset $gK_i$ specifies exactly the images of positions outside
$A_i$.  A uniform law on $G$ induces the uniform injection on those images.

The uniform marginal bounds above therefore give the hypothesis of
the simultaneous trimming lemma~\cite[Lemma~\ref*{l-l:trim}]{partialFourier} for every fixed equal partition.  For the averaged
estimate of Proposition~\ref{sweeps:averaged}, choose a uniform random
partition first: each complementary list is uniform, up to irrelevant
ordering, so the expected sum of its $q$ marginal TV errors is $o(1)$.
At least one deterministic partition has sum $o(1)$, and we fix it.
For clarity, the clipping is as follows.  Discard all cosets in any system
whose original mass exceeds twice their uniform mass.  The mass in such
cosets is at most twice the corresponding marginal TV error.  The union
therefore has mass $\delta=o(1)$.  Conditioning the original law $\mu$ on
its complement gives a probability $\nu$ with
\begin{equation}\label{sweeps:coset-bound}
 d_{\rm TV}(\mu,\nu)=\delta,\qquad
 \nu(gK_i)\le A/[G:K_i].
\end{equation}
One may take $A=3$ for sufficiently large $d$, or the convenient $A=4$.
The original permutation has unbiased sign, since toggling one fixed fair
switch toggles parity.  Thus
\begin{equation}\label{sweeps:sign}
 |\mathbb E_\nu\operatorname{sgn}|
 \le\delta/(1-\delta).
\end{equation}
This records the cap and retained mass needed for the applications below.

There is also a direct path-event construction of coset domination, which
uses the stronger conditional energy estimate rather than marginal TV.

\begin{proposition}[Trajectory-event truncation]\label{sweeps:path-truncation}
Use $q=16$ blocks, let $b$ be any exponent allowed by
Proposition~\ref{sweeps:two-sweep} for $\varepsilon=1/16$, and take an
integer $k$ with $bk\ge12$.  For the shuffle law $\mu$ after $T=2kd$ steps,
there is a probability $\nu$ and $\theta=1-\delta$ such that
\[
 \delta\le16n^{5-bk},\qquad d_{\rm TV}(\mu,\nu)\le\delta,
 \qquad (\nu*U_{K_i})(g)\le e/(\theta|G|)\quad(i,g).
\]
Moreover $|\mathbb E_\nu\operatorname{sgn}|\le\delta/\theta$.
\end{proposition}
\begin{proof}
For each complement fix an ordering $z_1,\ldots,z_{n-n/16}$.
For its $j$th card let $G_{i,j}$ be the event that the conditional final
energy given the preceding paths is at most $n^{-4}$.
The iterated energy bound and Markov's inequality give
$\Pr(G_{i,j}^c)\le n^{4-bk}$.  The event is measurable from those preceding
paths.  On it each conditional endpoint probability is at most
\[
 (n-j+1)^{-1}+n^{-2}\le\frac{1+1/n}{n-j+1}.
\]
Intersect all these events over the sixteen orderings, obtaining $G_*$ with
probability $\theta$ and the stated bound on $\delta$.

Fix one ordering and a tuple of distinct target endpoints.  Drop the events
from the other orderings.  Condition on the first $j-1$ full paths to bound
the last target endpoint on $G_{i,j}$; then drop $G_{i,j}$ and repeat
backwards.  Every conditioning uses the original shuffle law and an event
measurable from the preceding paths.  The probability of the target tuple
together with $G_*$ is consequently at most
\[
 \frac{(1+1/n)^{n-n/16}}{(n)_{n-n/16}}\le\frac e{(n)_{n-n/16}}.
\]
Let $\nu$ be the endpoint law conditional on the trajectory event $G_*$.  Right convolution by $U_{K_i}$ fills
each endpoint fiber uniformly, and each fiber has $(n/16)!$ elements.
This gives the density bound.  Conditioning on an event of probability
$\theta$ changes TV by at most $\delta$; the zero sign mean under $\mu$
gives the final estimate by splitting its expectation over $G_*$ and its
complement.
\end{proof}


\subsection{The representation inputs and their substitutions}\label{sweeps:transfer-inputs}
The input just obtained is a common coset cap for one probability law
$\nu$, together with a small sign coefficient.  The transfers we use have
different norm conclusions.  We state them explicitly so that the
following parameter choices can be checked without reconstructing the
companion's conventions.

For $u>0$, write
\[
 C_u=\sup_{a\ge1}\sum_{\tau\in\operatorname{Irr}(S_a)}(\dim\tau)^{-u}.
\]
The degree results in~\cite[Lemma~\ref*{l-l:degree-reciprocal-two}
and Theorem~\ref*{l-l:degree-all-powers}]{partialFourier} give $C_u<\infty$ and
\begin{equation}\label{sweeps:degree-input}
 \sum_{\lambda\vdash n,\,\lambda\notin\{(n),(1^n)\}}
                (\dim\lambda)^{-u}\longrightarrow0
 \qquad(n\longrightarrow\infty).
\end{equation}
Only the exponents $2,4,12,30,32$ will occur below.  The sharper
classical estimate of Liebeck and Shalev is, for each fixed $u>0$,
\begin{equation}\label{sweeps:LS}
 \sum_{\lambda\vdash a}(\dim\lambda)^{-u}=2+O(a^{-u})
 \qquad(a\longrightarrow\infty),
\end{equation}
by~\cite[Theorem~2.6]{liebeck-shalev2004}.

Suppose $K_1,\ldots,K_q$ permute disjoint equal blocks and
$\nu(gK_i)\le A/[G:K_i]$.  For a unitary irreducible $\rho$ of degree
$D$, use the mass Fourier transform
$\widehat\nu(\rho)=\sum_g\nu(g)\rho(g)$ and the unnormalized
Hilbert--Schmidt norm.  The orbit-span transfer gives
\begin{equation}\label{sweeps:op-bound}
 \|\widehat\nu(\rho)\|_{\rm op}
 \le\sqrt{AqC_u}\,D^{-1/2+(u+2)/(2q)}
 .
\end{equation}
This is Proposition~\ref*{l-l:orbit-span} of~\cite{partialFourier}.
The coefficient-evaluation transfer~\cite[Proposition~\ref*{l-l:coefficient-operator}]{partialFourier} gives the same bound
at $u=2$: its evaluation constant is the sum of squared degrees of
block types of degree at most $D^{1/q}$, which is at most
$C_2D^{4/q}$.  It counts each irreducible coefficient space once,
regardless of its multiplicity in $\rho$.
The pointwise-character and coset-Plancherel arguments instead give
\begin{align}
 \|\widehat\nu(\rho)\|_{\rm HS}^2
 &\le A^2qC_u D^{-1+(u+4)/q},\label{sweeps:pointwise-HS-eq}\\
 \|\widehat\nu(\rho)\|_{\rm HS}^2
 &\le AqC_u D^{-1+(u+2)/q},\label{sweeps:coset-HS-eq}
\end{align}
respectively~\cite[Propositions~\ref*{l-l:coarse-character-lift}
and~\ref*{l-l:isotypic}]{partialFourier}.
These are distinct proofs in the companion.  The second inequality
has stronger norm information than \eqref{sweeps:op-bound}; below we
also retain the applications of the operator argument itself.
For $A=4$, $u=2$ and large block size, the reciprocal-square sum for
each block is at most $3$, so the first Hilbert--Schmidt bound has
constant $48q$ and exponent $-1+6/q$.

\subsection{Parameter choices and full-deck conclusions}\label{sweeps:assembly}
All six applications now use probability laws. The ordinary marginal
bounds give such a law by the common coset deletion and normalization
in Section~\ref{sweeps:truncation}. The trajectory construction instead
conditions on its predictable path event. Both give a law $\nu$ within
$\delta=o(1)$ of the physical block law $\mu$, with sign coefficient
$o(1)$. Table~\ref{sweeps:transfer-table} records the resulting
substitutions in the three norm bounds above.

The integer conditions are part of the choices. In the trajectory row,
$b$ is any exponent allowed by Proposition~\ref{sweeps:two-sweep} for
$\varepsilon=1/16$, $k$ is a fixed integer with $bk\ge12$, and $p$ is any fixed integer satisfying
$3p/4-2\ge2$; in particular $p=6$ works. Its retained mass is
$\theta=1-\delta$. In the eight-orbit row, any fixed integer $r>300$
is allowed, and $r=400$ gives time $8000d$. In the available-count row,
choose an integer $b_0>2$ with $(1-1/512)^{b_0-2}\le2^{-6}$.

\begingroup
\small
\setlength{\tabcolsep}{4pt}
\renewcommand{\arraystretch}{1.2}
\begin{longtable}{@{}p{.33\linewidth}p{.29\linewidth}r r r@{}}
\caption{Sweep bounds and their Fourier applications. Here $q$ is the
number of label blocks, $A$ is the common coset cap, and $p$ is the
number of independent shuffle blocks. The exponent $\kappa$ is the
resulting nonsign Plancherel power.}\label{sweeps:transfer-table}\\
\toprule
Input and block parameters & Transfer and norm exponent
 & $p$ & $\kappa$ & Total time$/d$\\
\midrule
\endfirsthead
\toprule
Input and block parameters & Transfer and norm exponent
 & $p$ & $\kappa$ & Total time$/d$\\
\midrule
\endhead
\bottomrule
\endfoot
\phantomsection\label{sweeps:assembly-trajectory}
Trajectory truncation,\newline
Proposition~\ref{sweeps:path-truncation};\newline
$T/d=2k$, $q=16$, $A=e/\theta$
&
Coefficient evaluation,\newline
\eqref{sweeps:op-bound}, $u=2$;\newline
operator exponent $\alpha=3/8$
& $p$ & $2-3p/4$ & $2pk$\\
\phantomsection\label{sweeps:assembly-eight-orbits}
Two-sweep energy, \eqref{sweeps:eighth-bound};\newline
$T/d=2r$, $q=8$, $A=3$
&
Orbit, \eqref{sweeps:op-bound}, $u=2$;\newline
operator exponent $\alpha=1/4$
& $10$ & $-3$ & $20r$\\
\phantomsection\label{sweeps:assembly-blocker-orbits}
Refreshed blocker count,\newline
$\varepsilon=1/32$;\newline
$T/d=2048$, $q=32$, $A=4$
&
Orbit, \eqref{sweeps:op-bound}, $u=12$;\newline
operator exponent $\alpha=9/32$
& $32$ & $-16$ & $65536$\\
\phantomsection\label{sweeps:assembly-hole-orbits}
Refreshed available count,\newline
$\varepsilon=1/128$;\newline
$T/d=b_0$, $q=128$, $A=4$
&
Orbit, \eqref{sweeps:op-bound}, $u=30$;\newline
operator exponent $\alpha=3/8$
& $64$ & $-46$ & $64b_0$\\
\phantomsection\label{sweeps:assembly-colored-HS}
Uniform signed energy,\newline
\eqref{sweeps:colored-marginal};\newline
$\varepsilon=1/16$, $r=2048$,\newline
$T/d=2r$, $q=16$, $A=4$
&
Pointwise character,\newline
\eqref{sweeps:pointwise-HS-eq};\newline
$\|\widehat\nu\|_{\HS}^2\le48qD^{-5/8}$
& $8$ & $-4$ & $32768$\\
\phantomsection\label{sweeps:assembly-collision-HS}
Averaged collision bound,\newline
Proposition~\ref{sweeps:averaged};\newline
$q=128$, $\varepsilon=1/q$, $R=128q$,\newline
$T/d=2R$, $A=3$
&
Coset Plancherel,\newline
\eqref{sweeps:coset-HS-eq}, $u=32$;\newline
HS-squared exponent;\newline
$\beta=47/64$
& $64$ & $-46$ & $2097152$\\
\end{longtable}
\endgroup

For the first four rows, the norm bound has the form
$\|\widehat\nu\|_{\op}\le K D^{-\alpha}$; hence its nonsign
Plancherel summand after $p$ factors is at most
$K^{2p}D^{2-2p\alpha}$. For the last two rows the bound is
$\|\widehat\nu\|_{\HS}^2\le K D^{-\beta}$, giving
$K^pD^{1-p\beta}$. These are exactly the two substitutions in
Section~\ref{sec:probability-completion}, and give the displayed
$\kappa$. The trajectory cap $e/\theta$ is bounded as $\theta\to1$;
all other prefactors are fixed. The reciprocal-degree estimates
\eqref{sweeps:degree-input} at powers $2,2,12,30,4,32$, respectively,
make the six nonsign sums tend to zero. The pointwise-character row
also follows directly from the Liebeck--Shalev estimate
\eqref{sweeps:LS} at powers $2$ and $4$.

The probability completion now applies once to every row. The sign
bound is \eqref{sweeps:sign}, or $\delta/\theta$ for the trajectory
event. Replacing the $p$ factors costs at most $p\delta=o(1)$.
Since $T$ is divisible by $d$, the original product law is
$\mu^{*p}=q_d^{*(pT)}$ and has the physical time in the last column.
All constants are absolute and all conclusions hold for sufficiently large
integer $d$, uniformly over deterministic starting decks. There is also a
one-card lower bound: after $t<d$ layers a specified card has at most
$2^t\le n/2$ possible positions, so its variation distance from uniform
is at least $1-2^t/n\ge1/2$. Projection cannot increase variation distance, and
therefore the full-deck threshold-$1/4$ mixing time is at least $d$.
The common full-permutation support
lower bound $1-2^{tn/2}/n!$ tends to one for $t\le d$; hence each displayed
application gives $t_{\rm mix}(d)=\Theta(d)$ in physical-shuffle time.
